\subsection{局部化原理. 局部可忽略集}

命题 4（局部化原理）. --- 设 f 是 X 到拓扑空间 F 中的一个映射。假设对每个 \(x \in X\)，存在 x 的一个可积邻域 \(V_{x}\) 以及 X 到 F 中的一个可测映射 \(g_{x}\)，使得 \(f(y) = g_{x}(y)\) 在 \(V_{x}\) 中几乎处处成立。那么 f 是可测的。

设 K 是一个紧集；存在有限多个点 \(x_{i} \in K\)，使诸 \(V_{x_{i}}\) 构成 K 的一个覆盖。由此立得（§4, No.~9, 引理）：存在可忽略集 \(N \subset K\) 以及由可积集 \(M_{j}\) 组成的 K - N 的一个有限分割，使得每个集合 \(K \cap V_{x_{i}}\) 都是 N 的一个子集与若干个 \(M_{j}\) 的并，且在每个 \(M_{j}\) 上，f 等于某个函数 \(g_{x_{i}}\)。现在，对每个 \(M_{j}\)，存在可忽略集 \(N_{j} \subset M_{j}\) 以及 \(M_{j} - N_{j}\) 的由紧集序列 \(K_{nj}\)（ \(n \in N\) ）构成的一个分割；另一方面，对每个 \(K_{nj}\)，存在可忽略集 \(P_{nj} \subset K_{nj}\) 以及 \(K_{nj} - P_{nj}\) 的由紧集序列 \(K_{mnj}\)（ \(m \in N\) ）构成的一个分割，使得 f 到每个 \(K_{mnj}\) 上的限制是连续的。由于 N、诸 \(N_{j}\) 与诸 \(P_{nj}\) 的并是可忽略的，故 f 可测。

因此，可测函数的概念是一个局部性质的概念。

定义 3. --- 称集合 \(A \subset X\) 是（关于测度 \(\mu\) ）局部可忽略的，如果对每个 \(x \in X\)，存在 x 的一个邻域 V，使得 \(V \cap A\) 是可忽略的。

由局部化原理，每个局部可忽略集都是可测的。可忽略集的性质（ \(\S2\) ）表明：局部可忽略集的每个子集是局部可忽略的，且局部可忽略集的每个可数并是局部可忽略的。

命题 5. --- 要使集合 A 局部可忽略，必须且只需：对每个紧集 K，A∩K 是可忽略的。

由于 X 的每点都有紧邻域，条件显然是充分的。它是必要的，因为若对每个 \(x \in K\)，存在 x 的一个邻域 \(V_{x}\)，使得 \(A \cap V_{x}\) 是可忽略的，则存在有限多个点 \(x_{i} \in K\)，使诸 \(V_{x_{i}}\) 构成 K 的一个覆盖，于是 \(A \cap K\) 含于诸可忽略集 \(A \cap V_{x_{i}}\) 的并。

推论 1. --- 要使集合 A 可忽略，必须且只需它是局部可忽略的且外测度有限。

条件显然是必要的。反之，若条件满足，则 A 含于一个可积开集 G，而 G 是一个可忽略集 N 与紧集序列 \((\mathrm{K}_{n})\) 的并（ \(\S4\) , No.~6, 定理 4 的推论 2）；由于 \(A \cap N\) 与诸集合 \(A \cap K_{n}\) 都是可忽略的，它们的并 A 也是可忽略的。

推论 2. --- 每个局部可忽略的开集都是可忽略的（因而含于 \(\mu\) 的支集的余集中）。

因为开集 G 的外测度是诸测度 \(|\mu|(\mathrm{K})\)（对应于紧集 \(K \subset G\) 者）的上确界（ \(\S4\) , No.~6, 定理 4 的推论 4）；若 G 局部可忽略，则对含于 G 的每个紧集 K，有 \(|\mu|(\mathrm{K}) = 0\)，因此 \(|\mu|^{*}(\mathrm{G}) = 0\)。

推论 3. --- 在一个无穷远处可数的局部紧空间 X 中，每个局部可忽略集都是可忽略的。

由于 X 是紧集序列 \((\mathrm{K}_{n})\) 的并，每个局部可忽略集 A 都是诸可忽略集 \(A\cap K_{n}\) 的并，从而是可忽略的。

可以举出并非无穷远处可数的局部紧空间的例子，以及这种空间 X 上的测度的例子，使得 X 中存在局部可忽略但并非可忽略的集合（§1, 习题 5）。

推论 4. --- 设 f 是 X 到拓扑空间 F 中的一个映射。若 f 的不连续点所成的集合 N 是局部可忽略的，则 f 是可测的。

对每个紧集 \(K \subset X\)，\(K \cap N\) 是可忽略的（命题 5），因此对每个 \(\varepsilon > 0\)，存在紧集 \(K_{1} \subset K - (K \cap N)\)，使得 \(|\mu|(\mathrm{K} - \mathrm{K}_{1}) \leqslant \varepsilon\)（§4, No.~6, 定理 4），且由假设，f 到 \(K_{1}\) 上的限制是连续的，由此即得结论。

若 \(P\{x\}\) 是一个性质，则性质 \(\langle P\{x\}\) 局部几乎处处（关于 \(\mu\) ）成立」按定义等价于性质 \(\langle\) 使 \((x \in X\) 且非 \(P\{x\}\) ) 的 x 所成之集是（关于测度 \(\mu\) ）局部可忽略的」。若 F 是任意一个集合，则关系 \(\langle f(x) = g(x)\) 局部几乎处处成立」在 X 到 F 中的映射所成之集上是一个等价关系。特别地，若 F 是一个向量空间，则称 X 到 F 中使 \(f(x) = 0\) 局部几乎处处成立的映射 f 为局部可忽略的。我们把这些概念中与 §2 的 No.~4、5、6 里对几乎处处相等的函数所列举的性质相对应的大多数性质，留给读者去建立。我们只限于指出：若 X 到 Hausdorff 拓扑空间 F 中的两个连续映射 \(f, g\) 局部几乎处处相等，则依据命题 5 的推论 2，它们几乎处处相等（因而在 \(\mu\) 的支集的每一点上相等（§2, No.~4, 命题 9））；不过，我们明确陈述下面的命题，它是局部化原理的一个直接推论：

命题 6. --- 设 f 是 X 到拓扑空间 F 中的一个可测映射。每个与 f 局部几乎处处相等的 X 到 F 中的映射都是可测的。

